\section{A cubic bound for the exceptional classes of a Mumford square}
\label{sec:mumford}
This section treats a family whose exceptional classes are not Weil
classes. Its purpose is to reduce the two-dimensional space of classes
to one detecting subvariety, with a degree estimate.

Mumford's construction gives Shimura curves of polarized abelian
fourfolds whose general member has endomorphism algebra $\Q$, although
its Hodge group is smaller than the full symplectic group; see
\cite[\S1.1]{Noot}. Galluzzi describes Hazama's exceptional classes on
their squares in \cite[Introduction, p.~727;
Propositions~4.1--4.2, pp.~733--734]{Galluzzi}.
Moonen--Zarhin \cite[Theorem~0.1(iii), p.~2]{MoonenZarhin}
identify these as exceptional classes which are not of Weil type.
The representation calculation below recalls that description. The
additional conclusion is the cubic estimate obtained by composing one
algebraic correspondence with itself.

\subsection{The exceptional plane}

Let $(A,\theta)$ be a polarized abelian fourfold whose Hodge group is the
full group of a Mumford family. Thus
\[
 G^{\mathrm{sc}}_{\C}\simeq\mathrm{SL}_2^3,\qquad
 V_{\C}:=H^1(A,\C)\simeq U_1\otimes U_2\otimes U_3,
 \qquad \dim U_i=2.
\]
Here $G^{\mathrm{sc}}$ is the simply connected cover of $G$.
The three factors descend by the transitive action associated with a
totally real cubic field $F$. The hypothesis concerns the Hodge group
of this fibre, not merely its membership in the family. Put
\[
 P=H^2(A,\Q)_{\mathrm{prim}},\qquad
 Q_2(x,y)=\int_Axy\theta^2,\qquad
 \Theta=\theta_1+\theta_2\text{ on }A^2,\qquad
 \mu=\int_A\theta^4/4!.
\]

\begin{lemma}\label{lem:mumford-plane}
The primitive space $P$ has dimension $27$ and
$\operatorname{End}_{\mathrm{HS}}(P)=F$. The rational Hodge classes
of codimension two on $A^2$ form the direct sum of the six-dimensional
space of divisor products and a two-dimensional space
\[
 E\simeq F^0:=\ker(\operatorname{Tr}_{F/\Q}:F\longrightarrow\Q).
\]
The identification is induced by
$P\otimes P\simeq\operatorname{End}_{\Q}(P)$ using $Q_2$.
\end{lemma}

\begin{proof}
The transposition sign in $V^{\otimes2}$ gives
\[
 \bigwedge\nolimits^2V_{\C}
  =(\mathbf1,\mathbf1,\mathbf1)
   \oplus(\mathbf1,\mathbf3,\mathbf3)
   \oplus(\mathbf3,\mathbf1,\mathbf3)
   \oplus(\mathbf3,\mathbf3,\mathbf1).
\]
The first summand is the polarization. The other three are
nonisomorphic absolutely irreducible representations of dimension nine.
Their transitive descent has commutant the defining cubic field $F$.
Since the Hodge group is $G$, this commutant is exactly the Hodge
endomorphism algebra of $P$.

For completeness, the invariant dimensions in the five K\"unneth
components of $H^4(A^2)$ are $1,1,4,1,1$. One can compute the two outer
dimensions without decomposing all exterior powers. The
$\mathrm{SL}_2$-invariants in $U^{\otimes4}$ have the $S_4$ character
$(2,0,2,-1,0)$, on the classes of cycle types
$1^4,\,2\,1^2,\,2^2,\,3\,1,\,4$. They are spanned by the three pairings
of four vectors, with their single alternating relation. Cubing the
character for the three factors gives $(8,0,8,-1,0)$.
Its sign multiplicity is
$(8+3\cdot8-8)/24=1$, and its sign multiplicity after restriction
to $S_3$ is $(8-2)/6=1$. These are respectively the invariants in
$\bigwedge^4V$ and in $V\otimes\bigwedge^3V$. The middle invariant
dimension is $\dim\operatorname{End}_G(\bigwedge^2V)=4$.

The three independent divisor classes on $A^2$ are
$\theta_1,\theta_2$ and the mixed Poincar\'e class $\ell$.
Their six quadratic products are independent. Four occupy the four
outer K\"unneth components; in the middle component,
$\theta_1\theta_2$ and $\ell^2$ are independent since $\ell^2$ has a
nonzero primitive coevaluation component. Under the identification of
the primitive invariant tensors with $F$, their middle primitive
line is $\Q\operatorname{id}_P$. Its complement is $F^0$.
\end{proof}

Let $\delta_{\mathrm{exc}}$ be the least $\Theta$-degree of an
irreducible codimension-two subvariety of $A^2$ whose projection to $E$
is nonzero. Such subvarieties have dimension six. The projection is the
one furnished by the decomposition in Lemma~\ref{lem:mumford-plane}.
As before the value is $+\infty$ if no detector exists.

\begin{theorem}\label{thm:mumford-cubic}
For every fixed nonzero rational class $t\in E$, there is a constant
$C_t>0$ such that
\[
 \delta_{\mathrm{exc}}\leq c_t\leq C_t\delta_{\mathrm{exc}}^3.
\]
The constants may be chosen uniformly on the Hodge-generic fibres of
one fixed polarized Mumford family, after a fixed finite cover making
$E$ and $t$ global. In particular, one nonzero algebraic exceptional
class implies algebraicity of the whole exceptional plane at a
Hodge-generic fibre.
\end{theorem}

The degree exponent has a simple explanation. An algebraic detector
gives a trace-zero element $a$ of a cubic field at linear cost.
Composing its correspondence produces $a^2$ at quadratic cost.
The trace-zero parts of $a$ and $a^2$ are independent. The area of the
lattice they span, and the cost of clearing the coefficients of a
fixed class, are cubic in the original degree.

\subsection{Algebraic operations and their costs}

We give the operators explicitly, to avoid assuming an algebraic
projector onto a Hodge endomorphism field.
For an integer $m\geq2$, the polynomial
\[
 \pi_k=
 \prod_{\substack{0\leq j\leq8\\j\ne k}}
 \frac{[m]^*-m^j}{m^k-m^j}
\]
in algebraic correspondences is the cohomological projector onto
$H^k(A)$. Let $L$ denote multiplication by $\theta$, and let $*$ denote
Pontryagin product. The lowering operator is
\[
 \Lambda(x)=\frac{\theta^3}{6\mu}*x.
\]
The identity follows by writing the polarization in a symplectic basis:
Pontryagin product contracts one paired pair of basis covectors, with
coefficient $\mu$, cancelled by the displayed normalization.
It is the usual cohomological Lefschetz operator, including for a
nonprincipal polarization. In particular on $H^2(A)$,
\begin{equation}\label{eq:mumford-operators}
 \pi_\theta=\tfrac14L\Lambda,\qquad
 \pi_P=1-\pi_\theta,\qquad
 R=(L^2)^{-1}
   =\left(\tfrac14\pi_P+\tfrac1{36}\pi_\theta\right)\Lambda^2.
\end{equation}
Indeed $\Lambda^2L^2$ is $4$ on primitive degree two and $36$ on
$\Q\theta$. These also follow from
$\Lambda L^jx=j(4-k-j+1)L^{j-1}x$ for a primitive degree-$k$ class.
All the operators in \eqref{eq:mumford-operators} are algebraic.

The identity tensor $I_P\in P\otimes P$ is algebraic: take the primitive
$H^6\otimes H^2$ component of the diagonal and apply $R$ to its first
factor. For a class $z\in P\otimes P$ write
\[
 a=\Phi(z),\qquad
 \Phi\left(\sum_i u_i\otimes v_i\right)(x)
   =\sum_i Q_2(x,u_i)v_i.
\]
Its rational trace is the algebraic intersection
\[
 \operatorname{Tr}_{P}(a)
   =\int_{A^2}zI_P\theta_1^2\theta_2^2.
\]
Thus the operation
\begin{equation}\label{eq:mumford-trace}
 z\longmapsto z-\frac{\operatorname{Tr}_{P}(a)}{27}I_P
\end{equation}
is algebraic with a fixed denominator. At a Hodge-generic fibre
$a\in F$ and $\operatorname{Tr}_{P}(a)=9\operatorname{Tr}_{F/\Q}(a)$.
Projecting a cycle to the $(2,2)$ K\"unneth component, then to
$P\otimes P$, and then subtracting this trace is exactly its
exceptional projection.

\begin{lemma}\label{lem:mumford-cost}
Give a rational cycle $Z=\sum_i r_iZ_i$ the signed cost
$|Z|_\Theta=\sum_i|r_i|\deg_\Theta Z_i$.
The fixed operators just described change cost $B$ to at most $CB$,
with fixed denominators. Composition of two correspondences on $A$
of costs $B_1,B_2$ has a signed representative of cost at most
$C'B_1B_2$. The constants depend only on the fixed data and are uniform
in the polarized family.
\end{lemma}

\begin{proof}
Pullbacks and pushforwards by multiplication maps, products with fixed
divisors, and Pontryagin products with fixed cycles have linear cost
bounds. Translate cycles generally before an intersection. Translation
preserves their classes and degrees, and the resulting proper
intersection of effective cycles is effective.
Fixed rational coefficients and the denominators of the displayed
operators are harmless.

Here is an intrinsic estimate covering both these assertions and
composition. On a polarized abelian variety use the flat K\"ahler
metric whose form represents the polarization. An effective
$d$-cycle $W$ has mass $\deg W/d!$. A harmonic form is constant, so
its pointwise norm is its $L^2$ norm divided by
$\sqrt{\operatorname{Vol}}$. Comass is at most pointwise norm.
Duality against all harmonic forms of unit norm consequently gives
\begin{equation}\label{eq:mumford-mass}
 \|[W]\|_{L^2}\leq
       \frac{\deg W}{d!\sqrt{\operatorname{Vol}}}.
\end{equation}
Exterior multiplication between two fixed Euclidean exterior powers
is a bounded bilinear map, with a constant depending only on their
degrees and the dimension.

For effective codimension-four correspondences $C,D$ on $A^2$, of
degrees $b,c$, intersect $p_{12}^*C$ and $p_{23}^*D$ on $A^3$
after translating $D$ generally. The action of translations on
the two coordinates of $D$ is transitive on $A^2$, so generic
transversality makes this intersection proper.
Its pushforward by $p_{13}$ represents the composition.
With $\Theta_3=\theta_1+\theta_2+\theta_3$, the pullback degrees are
\[
 \deg_{\Theta_3}p_{12}^*C
   =\binom84(4!\mu)b,\qquad
 \deg_{\Theta_3}p_{23}^*D
   =\binom84(4!\mu)c.
\]
Apply \eqref{eq:mumford-mass}, the bilinear exterior estimate and pairing with
$\Theta_3^4$. The degree of the intersection is at most a fixed
multiple of $bc$. Since
$p_{13}^*\Theta\leq_{\mathrm{nef}}\Theta_3$, the degree of its
pushforward is no greater. Expand signed cycles linearly to obtain
the general assertion. The flat metric estimate depends only on
dimension and polarization volume; the fixed maps have the same
norms and degrees throughout the family.
\end{proof}

\subsection{Proof of the cubic estimate}

\begin{proof}[Proof of Theorem~\ref{thm:mumford-cubic}]
Let $Y$ be a detector of degree $D$. Fix the order
\[
 \mathcal O=F\cap\operatorname{End}_{\Z}(P_{\Z})
\]
for the fixed primitive cohomology lattice. The fixed projectors and
trace subtraction, followed by one fixed clearing of denominators,
give a nonzero $\alpha\in\mathcal O^0$ represented by an integral
signed codimension-two cycle of cost at most $C_1D$.

Multiply its tensor by $\theta_1^2$ to make a codimension-four
correspondence on $A$ acting as $\alpha$ on $P$; use the fixed
projectors to remove its other cohomological components.
Compose it with itself and use $R$ on the degree-six factor to
return to codimension two. Lemma~\ref{lem:mumford-cost} gives a
representative of $\alpha^2$ of cost $O(D^2)$ and a fixed
denominator. Trace subtraction and multiplication by one fixed
positive integer $r$ give
\[
 \beta=r\bigl(3\alpha^2-
        \operatorname{Tr}_{F/\Q}(\alpha^2)\operatorname{id}\bigr)
      \in\mathcal O^0
\]
with an integral signed representative of cost at most $C_2D^2$.
The trace term has this cost bound by the estimate below.
A nonzero trace-zero element of a cubic field is nonrational and
generates the field. If $\beta$ were a rational multiple of $\alpha$,
a polynomial of degree two would annihilate $\alpha$, which is
impossible. Thus $\alpha,\beta$ span $F^0$.

Give $F^0_{\R}$ the norm
$\|a\|_F^2=\operatorname{Tr}_{F/\Q}(a^2)$.
The adjoint on $F$ is the identity, since the three scalar summands
of $P_\C$ are mutually orthogonal. If $z_a$ is the tensor
corresponding to $a$, then
\[
 \int_{A^2}z_a^2\Theta^4
   =6\operatorname{Tr}_{P}(a^2)
   =54\operatorname{Tr}_{F/\Q}(a^2).
\]
Each factor of $z_a$ is primitive of degree two; hence
$\Theta^5z_a=0$. The primitive Hodge-star formula gives
$\|z_a\|_{L^2}^2=\frac94\|a\|_F^2$.
Since $\operatorname{Vol}(A^2,\Theta)=\mu^2$, inequality
\eqref{eq:mumford-mass} for signed six-cycles of cost $B$ yields
\[
 \|a\|_F\leq\frac{2B}{3\cdot6!\mu}.
\]
In particular $\|\alpha\|_F=O(D)$,
$\operatorname{Tr}_{F/\Q}(\alpha^2)=O(D^2)$, and
$\|\beta\|_F=O(D^2)$.

Choose $\delta\geq1$ with $\delta t\in\mathcal O^0$ and set
$h=[\mathcal O^0:\Z\alpha+\Z\beta]$.
The positive integral trace pairing on $\mathcal O^0$ has
covolume at least one. The area formula therefore gives
$h\leq\|\alpha\|_F\|\beta\|_F=O(D^3)$.
There are integers $u,v$ such that
\[
 h\delta t=u\alpha+v\beta.
\]
The two-dimensional determinant formula bounds
$|u|\leq\delta\|t\|_F\|\beta\|_F=O_t(D^2)$ and
$|v|\leq\delta\|t\|_F\|\alpha\|_F=O_t(D)$.
The multiplier $h\delta$ is consequently $O_t(D^3)$, and the
signed representative on the right has total cost
$|u|C_1D+|v|C_2D^2=O_t(D^3)$.
Separate its positive and negative coefficients to obtain the
upper bound for $c_t$.

Conversely, every effective difference representing a nonzero
multiple of $t$ has an irreducible component with nonzero
exceptional projection. Its degree is no greater than the
corresponding numerator degree. This proves the lower bound.
The order, its trace pairing, the fixed operators and the norm of
$t$ are fixed in the family, proving the asserted uniformity.
\end{proof}

\begin{corollary}\label{cor:mumford-family}
On a fixed connected Mumford Shimura curve, suppose that detectors
for its exceptional plane have degrees bounded by one integer $D$
at infinitely many CM points. Then, for every fixed global
$t\in E$, the complexity $c_t$ is bounded on every fibre.
\end{corollary}

\begin{proof}
The locus with a degree-$D$ detector is closed: use the proper
bounded-cycle parameter spaces of Lemma~\ref{lem:parameters}, and
select components on which the projection to the global flat plane
is nonzero. This uses a cohomological projection to specify a
locus, not an algebraic correspondence projecting to $F$.
An infinite subset of the irreducible curve is Zariski dense, so
every fibre has a degree-$D$ detector. Apply
Theorem~\ref{thm:mumford-cubic} on its dense Hodge-generic locus.
The resulting uniform bound for $c_t$ extends to every fibre by
the closedness of Lemma~\ref{lem:closed}.
\end{proof}

The Hodge-generic qualification in the pointwise theorem is necessary
for this proof. At a CM fibre,
$\operatorname{End}_{\mathrm{HS}}(P)$ may be larger than $F$.
The algebraic trace subtraction above then projects to the
trace-zero part of that larger algebra; it need not extract $F^0$
from an arbitrary detecting cycle. If a class already lies in
$F^0$, the field-composition argument applies. Corollary
\ref{cor:mumford-family} passes through the generic fibres and does
not assume a pointwise extraction operator at CM fibres.

